\chapter{CJK Scripts}
\label{ch:cjk}

OmniLaTeX provides first-class support for Chinese, Japanese, and Korean (CJK)
scripts through \pkg{luatexja} and the Noto CJK font family.  When you select
a CJK language (e.g.\ \texttt{language=simplifiedchinese}), OmniLaTeX
automatically configures the appropriate fonts and script settings so you can
mix CJK and Latin text seamlessly.

\section{Chinese (Simplified)}

Simplified Chinese uses the Noto Sans CJK SC and Noto Serif CJK SC font
families, loaded through \pkg{luatexja}'s font mechanism.

\begin{verbatim}
\setCJKMainFont{Noto Serif CJK SC}
\setCJKSansFont{Noto Sans CJK SC}
\setCJKMonoFont{Noto Sans Mono CJK SC}
\end{verbatim}

OmniLaTeX calls \verb|\setCJKScript{sc}| internally whenever
\texttt{language=simplifiedchinese} is detected, which tells \pkg{luatexja}
to activate the Simplified Chinese OpenType script features (GB-standard
glyphs, punctuation rules, and proportional spacing).

A complete example lives in \texttt{examples/cjk-chinese/}, which
demonstrates:

\begin{itemize}
  \item Mixed Latin--Chinese paragraphs with automatic inter-script spacing.
  \item Serif and sans-serif Chinese body text.
  \item Chinese-style punctuation (full-width commas, periods, enumeration
        commas ``\punct{、}'').
  \item Headings, lists, and captions rendered in Chinese.
\end{itemize}

To compile a Simplified Chinese document:

\begin{verbatim}
\documentclass{omnilatex}
\usepackage[language=simplifiedchinese]{omnilatex-lang}
\begin{document}
  \chapter{中文标题}
  这是一段简体中文示例。
\end{document}
\end{verbatim}

\section{Chinese (Traditional)}

Traditional Chinese switches to the TC (Traditional Chinese) variant of the
Noto CJK family:

\begin{verbatim}
\setCJKMainFont{Noto Serif CJK TC}
\setCJKSansFont{Noto Sans CJK TC}
\end{verbatim}

The script tag is \verb|\setCJKScript{tc}|, activated by
\texttt{language=chinese-traditional}.  This selects the correct glyph set
(e.g.\ 體 vs.\ 体, 裡 vs.\ 里) and adjusts punctuation rules for the
Traditional Chinese typographic convention used in Taiwan, Hong Kong, and
Macau.

\begin{verbatim}
\documentclass{omnilatex}
\usepackage[language=chinese-traditional]{omnilatex-lang}
\begin{document}
  \chapter{繁體中文標題}
  這是一段繁體中文範例。
\end{document}
\end{verbatim}

\section{Japanese}

Japanese documents use the Noto CJK JP fonts together with \pkg{luatexja}:

\begin{verbatim}
\setCJKMainFont{Noto Serif CJK JP}
\setCJKSansFont{Noto Sans CJK JP}
\setCJKMonoFont{Noto Sans Mono CJK JP}
\end{verbatim}

The script tag is \verb|\setCJKScript{jp}|, activated by
\texttt{language=japanese}.  In addition to standard CJK support,
OmniLaTeX enables \emph{furigana} (reading aids) for Japanese text via
\pkg{luatexja-ruby}.

A complete example lives in \texttt{examples/cjk-japanese/}, demonstrating
furigana placement, mixed-script paragraphs, and Japanese-style punctuation
and line-breaking.

\begin{verbatim}
\documentclass{omnilatex}
\usepackage[language=japanese]{omnilatex-lang}
\begin{document}
  \chapter{日本語の見出し}
  これは日本語の例です。
\end{document}
\end{verbatim}

\section{Korean}

Korean documents use the Noto CJK KR font family:

\begin{verbatim}
\setCJKMainFont{Noto Serif CJK KR}
\setCJKSansFont{Noto Sans CJK KR}
\end{verbatim}

The script tag is \verb|\setCJKScript{kr}|, activated by
\texttt{language=korean}.  Hangul syllables are rendered with proper
morphological spacing, and Korean-style punctuation (middle dots ``·'')
is supported.

A complete example lives in \texttt{examples/cjk-korean/}.

\begin{verbatim}
\documentclass{omnilatex}
\usepackage[language=korean]{omnilatex-lang}
\begin{document}
  \chapter{한국어 제목}
  이것은 한국어 예시입니다.
\end{document}
\end{verbatim}

\section{Auto-Detection}

When you pass a CJK language option to \pkg{omnilatex-lang}, OmniLaTeX
performs the following steps automatically:

\begin{enumerate}
  \item Detects the language family (SC, TC, JP, KR) from the option name.
  \item Calls \verb|\setCJKScript{<tag>}| to activate the correct OpenType
        script.
  \item Loads the matching Noto CJK font family for main, sans, and mono.
  \item Sets \pkg{luatexja} parameters for inter-character spacing and
        line-breaking rules appropriate to the script.
  \item Configures translation keys (chapter, figure, table, etc.) in the
        target language.
\end{enumerate}

You can override any of these defaults by setting CJK fonts or scripts
manually in your preamble \emph{after} loading \pkg{omnilatex-lang}.

\section{Ruby Annotations}

Ruby annotations (small phonetic glosses above or beside characters) are
supported via \pkg{luatexja-ruby}:

\begin{itemize}
  \item \verb|\ruby{漢字}{かんじ}| -- Japanese furigana (above).
  \item \verb|\furigana{漢字}{かんじ}| -- Alias with identical behaviour.
  \item \verb|\pinyin{汉字}{hànzì}| -- Chinese pinyin (above).
\end{itemize}

Control the ruby annotation size with \verb|\setRubyScale{0.5}| (default
\texttt{0.5}), which scales the annotation font relative to the base font
size.

\begin{verbatim}
\ruby{漢字}{かんじ}   % furigana
\pinyin{汉字}{hànzì}  % pinyin
\setRubyScale{0.45}    % smaller annotations
\end{verbatim}

\section{Vertical Writing Mode}

OmniLaTeX supports vertical (top-to-bottom, right-to-left) writing via the
\texttt{vertical} environment provided by \pkg{luatexja}:

\begin{verbatim}
\begin{vertical}
  この文章は縦書きです。右から左へ読んでください。
\end{vertical}
\end{verbatim}

For inline vertical text, use \verb|\verticaltext{縦書き}|.  Page layout in
vertical mode is adjusted automatically--page numbers and running headers
rotate to match the text direction.

\section{CJK--Latin Spacing}

When mixing CJK and Latin scripts, proper inter-character spacing prevents
the text from looking cramped.  OmniLaTeX configures \pkg{luatexja} to
insert a thin space (typically $\tfrac{1}{4}$~em) between CJK and Latin
characters by default.

You can fine-tune this behaviour:

\begin{verbatim}
\cjklatinspacing          % enable (default for CJK languages)
\nocjklatinspacing        % disable
\setCJKLatinSpacing{0.25} % set gap to 0.25 em
\end{verbatim}

The spacing value is a fraction of one em in the current font size.
